\documentclass[11pt]{article}
    \usepackage[utf8]{inputenc}
    \usepackage[T1]{fontenc}
    \usepackage[margin=1in]{geometry}
    \usepackage{hyperref}
    \usepackage{enumitem}
    \title{Coding a Lego Zoo}
    \author{Piter Garcia}
    \date{March 20, 2026}
    \begin{document}
    \maketitle
    \section*{Quick Scan}
    \begin{itemize}[leftmargin=*]
    \item Grade: Grade 1
    \item Workbook sessions: Session 13
    \item Class blocks: Day 3 12:25-1:15 Gargana, Day 3 2:15-3:05 Polfleit, Day 5 12:25-1:15 Montgomery
    \item Reference file: coding-a-lego-zoo.bib
    \end{itemize}
    \section*{School Planning Goals and Inclusion}
    \begin{itemize}[leftmargin=*]
    \item What is expected: Legos | Maze Mat | Coding cards | Mini toy animals | Dry erase markers
    \item What we achieved: This lesson points to local transcript or evidence files in the workspace so the packet can be revised from what actually happened in class.
    \item What needs improvement: stronger proof, cleaner assessment notes, and tighter lesson artifacts.
    \item How to improve: add transcript proof, one saved artifact, and clearer success criteria.
    \item Inclusion focus: use visual modeling, chunked directions, predictable routines, and flexible participation roles.
    \end{itemize}
    \section*{Official References}
    \begin{itemize}[leftmargin=*]
    \item \url{https://csteachers.org/k12standards/interactive/}
\item \url{https://assets.education.lego.com/v3/assets/blt293eea581807678a/blt593e0d51d19eceb5/62f4d13ce894104b27259cb8/Coding_express_teacher_guide_110822.pdf?locale=ja-jp}
\item \url{https://code.org/curriculum/elementary-school}

    \end{itemize}
    \end{document}
